problem:  
**Conjecture (Periodic Cohomology for Almost‑Maximal Torus Symmetry in Positive Curvature).**  
Let \( M^n \) be a compact, simply connected Riemannian manifold with strictly positive sectional curvature, and suppose it admits an effective isometric torus action of rank  
\[
r = \left\lfloor \frac{n+1}{2} \right\rfloor - 1,
\]  
that is, one below the maximal symmetry rank allowed by the Grove–Searle theorem.  
Assume further that all singular orbits of the \(T^r\)-action have codimension at most \(4\).

**Problem:**  
Must \(M\) have **4‑periodic rational cohomology**, meaning there exists a class \(x \in H^4(M;\mathbb{Q})\) such that the cup‑product map  
\[
H^i(M;\mathbb{Q}) \xrightarrow{\smile\, x} H^{i+4}(M;\mathbb{Q})
\]  
is an isomorphism for all admissible degrees \(i\)?  
In particular, does every such \(M\) share this characteristic cohomological periodicity property with the compact rank‑one symmetric spaces?

---

Why is it a "good" problem:  
This conjecture isolates the *first unexplored layer* of the Grove symmetry program beyond the classical maximal‑rank case. The Grove–Searle theorem completely classifies positively curved manifolds when the torus symmetry rank is maximal; what remains largely unknown is what happens at one step below this upper bound.

The problem’s simplicity hides its depth: it asks whether **large but not maximal symmetry** already enforces **periodicity in rational cohomology**, a hallmark shared by CROSSes. The codimension‑4 hypothesis reflects the natural geometric threshold where Wilking’s connectedness lemmas can produce 4‑periodicity in known cases—pushing just beyond the cases now controlled by existing theorems due to Wilking, Amann–Kennard, and Dessai.  

Resolving the conjecture would integrate techniques from **Riemannian geometry** (positive curvature and orbit geometry), **transformation group theory** (action stratification and isotropy analysis), and **rational homotopy theory** (periodicity and cohomological generation).  
A proof would establish a sharp new rigidity layer in positive curvature; a counterexample would reveal a previously unseen interaction between curvature and symmetry. The statement is crisp, geometrically natural, and precisely at a known boundary of current understanding—making it a “good” research‑level problem.